% 本文件由 LaTeX 排版 Agent 根据有效 translation.json 自动生成。
% author=Athanasius; work_id=2803; title=亚略的罢免

\chapter{亚略的罢免}

% structure: p0001 source=h1 target=omitted reason=page-title-already-rendered-by-parent

\Emph{亚历山大对亚略及其同党的罢免，以及关于此事的通谕信。}\par

\Person{亚历山大}，同他所爱的弟兄们，即\Place{亚历山大城}及\Place{马雷奥蒂斯}的司铎与执事们，聚集在一起，在主内向各位致意。\par

虽然你们已经在我写给\Person{亚略}及其同党的信中签字确认，劝诫他们弃绝其不虔敬，并服从纯正的天主教信德，也显示了你们的正见及对天主教会教义的一致性；然而，由于我又就\Person{亚略}及其同党的事写信给我们在各处的同工们，尤其因为你们当中有些人，诸如司铎\Person{卡雷斯}与\Person{皮斯图斯}，执事\Person{塞拉皮翁}、\Person{帕拉蒙}、\Person{佐西穆斯}和\Person{依勒内}，已加入\Person{亚略}及其同党，并甘愿与之同受罢免，我认为有必要召集你们这城里的圣职人员，并召来\Place{马雷奥蒂斯}地区的圣职人员，以便你们能够了解我现在所写的事，并证实你们的赞同，且一致同意罢免\Person{亚略}、\Person{皮斯图斯}及其同党。因为我所写的事，理当让你们知晓，且你们每人当衷心接纳，如同自己写的一般。\par

\Emph{抄本。}\par

致各地最亲爱的、极可敬的天主教会同工们，\Person{亚历山大}在主内致安。\par

1. 正如天主教会只有一个身体，且圣经中吩咐我们要保守合一与和平的纽带，故此，我们应当彼此通信，告知各自所行的事；这样，若一个肢体受苦，所有肢体都一同受苦；若一个肢体得荣耀，所有肢体都一同欢乐。如今，在这教区，此时有一些不法之人出去，他们是基督的仇敌，教导一种背教之事，人们有理由怀疑并称之为敌基督的先驱。我本欲对此事置之不提，希望这邪恶或许只在其拥趸中消耗殆尽，而不致延及他处，污染纯朴者的耳朵。然而，眼见那如今在\Place{尼科米底亚}的\Person{欧瑟比乌斯}，自以为教会的治理权在他手中，因为他对\Place{贝里图斯}的离弃尚未遭到报应，那时他便垂涎\Place{尼科米底亚}人的教会；如今竟开始支持这些背教者，并擅自四处为他们写信，企图设法引诱一些无知的人归向这最卑鄙、敌基督的异端；故此，我既知晓法律上所记的，不得不打破沉默，向你们众位揭示出来，使你们明了这些背教者是谁，以及他们所持异端的诡辩，且若\Person{欧瑟比乌斯}写信给你们，你们不必理会他，因为他现今意欲借这些人重新显扬他那久藏的宿怨，他假装为他们写信，事实上显然是要图谋自己的利益。\par

2. 那些成了背教者的人如下：\Person{亚略}、\Person{阿基里斯}、\Person{艾塔勒斯}、\Person{卡尔颇内斯}、另一位\Person{亚略}，和\Person{萨尔玛特斯}，曾为司铎；\Person{欧佐伊乌斯}、\Person{路西乌斯}、\Person{犹利乌斯}、\Person{梅纳斯}、\Person{赫拉迪乌斯}和\Person{盖乌斯}，曾为执事；与他们一起的还有\Person{塞昆杜斯}和\Person{特奥纳斯}，曾被称为主教。他们所发明并传布的、相反圣经的新奇之说如下：—— 天主并非永远是圣父，曾有一段时间，天主不是圣父。天主的圣言并非永恒，而是源自虚无；因为那自有者天主，从无中造了他，他原本并不存在；因此曾有一段时间他并不存在；因为圣子是一个受造物和工程。他在\OriginalTerm{本质}{essence}上也不像圣父；他既不是圣父的真正而本然的圣言；也不是他的真智慧；他是受造物之一，被称为圣言和智慧，乃是误用词语；因为他自己乃是由天主本有的圣言和在天主内的智慧所产生的，天主不仅用这智慧创造了其他一切，也创造了他。因此，他本性上如所有理性受造物一样，易于变化和更动。圣言与圣父的本质无关，是外来的，与圣父分离。圣父也不能被圣子所描述，因为圣言不完全、不准确地认识圣父，也不能完全看见圣父。而且，圣子也不知道自己本质的真实状况；因为他为我们而被造，好使天主借着他如同工具一样创造我们；倘若天主不愿创造我们，他便不会存在。因此，当有人问他们，天主的圣言是否能像魔鬼那样变化，他们竟毫不畏惧地说可以；因为他是受造之物，其本性易于变化。\par

3. 当\Person{亚略}及其同伙作出这些断言，并厚颜无耻地公开宣认时，我们与\Place{埃及}和\Place{利比亚}的主教们聚集，人数将近一百，我们绝罚了他们及其追随者。但\Person{优西比乌斯}及其同党却接纳他们共融，渴望使谎言与真理混杂，不虔与虔敬混杂。但他们不能如此，因为真理必要得胜；既没有\BibleQuote{光明与黑暗的共融}，也没有\BibleQuote{基督与\Person{贝里雅耳}的和谐}。因为谁曾听说过这样的论调呢？谁现在听了不感到惊异，不掩住耳朵，以免被这些言语玷污呢？谁听了\Person{若望}的话，\BibleQuote{在起初已有圣言，}\BibleRef{若 1:1} 而不谴责这些人所说的\Quote{曾有一个祂不存在的时刻}呢？或者，谁在福音中听到\BibleQuote{独生子}和\BibleQuote{万物是藉着祂而造成的}，而不憎恶他们宣称祂是\Quote{受造物之一}呢？因为祂怎能成为自己造成的万物之一呢？又怎能是独生子，若按照他们的说法，祂被算作其余中的一位，既然祂自己也是受造物和工程？祂又怎能是\Quote{生于虚无}的呢？圣父说：\BibleQuote{我的心灵涌溢美辞}，又说：\BibleQuote{我在晓明之前，由母腹中生了祢}？再者，祂怎能\Quote{不同性体于圣父}呢？既然祂是圣父圆满的\BibleQuote{肖像}和\BibleQuote{光辉}\BibleRef{希 1:3}，且祂自己说：\BibleQuote{谁看见了我，就看见了父}？如果圣子是\OriginalTerm{圣言}{\GreekText{Λόγος}}和\OriginalTerm{智慧}{\GreekText{Σοφία}}，怎么会有\Quote{祂不存在的时刻}呢？这就如同他们说天主曾一度没有圣言和智慧。祂怎能\OriginalTerm{易变易改}{\GreekText{τρεπτὸν} \GreekText{καὶ} \GreekText{ἀλλοιωτόν}}呢？祂亲口说：\BibleQuote{我在父内，父在我内}，又说：\BibleQuote{我与父原是一体}；并由先知说：\BibleQuote{看，我就是我，决不改变}？因为，尽管这句话可以指圣父，但现在更贴切地可以用于圣言：即，虽然祂成了人，祂并没有改变；正如宗徒所说：\BibleQuote{耶稣基督昨天、今天、直到永远，常是一样的}。谁竟能说服他们说祂是为我们而受造的呢？\Person{保禄}却写道：\BibleQuote{万有都是靠祂，并藉着祂而存在的}？\par

4. 至于他们亵渎的主张，认为\Quote{圣子并不完全认识圣父}，我们不必感到奇怪；因为他们一旦决意与基督为敌，就甚至反驳祂明确的话语，因为祂说：\BibleQuote{正如父认识我一样，我也认识父}\BibleRef{若 10:15}。那么，如果父只是部分地认识子，显然子就不完全认识父；但若这样说是不允许的，而是父完全认识子，那么显然正如圣父认识自己的\OriginalTerm{圣言}{\GreekText{Λόγος}}，同样地，圣言也认识自己的父，祂正是这圣言。\par

5. 我们常用这些论证和引用\WorkTitle{圣经}来驳倒他们；但他们却像变色龙一样变化，一再改变他们的立场，努力使自己应验那句话：\BibleQuote{恶人一堕入罪恶的深处，便肆无忌惮}。在他们之前已有许多异端，这些异端因为越过了应有的界限而陷入愚妄；但这些人费尽心机，用各种诡辩企图推翻圣言的天主性，相比之下，反而使那些异端显得有理，因为他们更接近\OriginalTerm{敌基督}{\GreekText{ἀντίχριστος}}。因此他们已被教会开除教籍并绝罚。我们为他们的丧亡而悲伤，尤其因为他们曾受过教会教义的教导，如今却背离了。然而我们并不感到非常惊奇，因为\Person{依默纳约}和\Person{菲肋托}也曾如此\BibleRef{弟后 2:17}，在他们之前有\Person{犹达斯}，他跟随了救主，后来却成为叛徒和背教者。关于这些事，我们并非没有受过教导；因为我们的主早已警告过我们：\BibleQuote{你们要谨慎，不要被人欺骗；因为将有许多人，假冒我的名字来说：我就是基督，时期近了；你们切不可跟随他们}\BibleRef{路 21:8}；而由我们救主亲自教导的\Person{保禄}则写道：\BibleQuote{在末后的时期，有些人要离弃信德，听从欺诈的神和魔鬼的训诲，这训诲是弃绝真理的}。\par

6. 既然我们的主和救主耶稣基督亲口教导了我们，又通过宗徒向我们指明了关于这类人的事，我们作为他们不虔行为的亲自见证者，正如我们所说，已经绝罚了所有此类人，并声明他们与公教信仰和天主教会格格不入。我们已将此事告知你们的虔诚，至爱和最受尊敬的同工，以便若他们有胆量到你们那里去，你们不要接纳他们，也不要顺应\Person{优西比乌斯}或任何其他为他们写信之人的意愿。因为我们基督徒应当远离一切说或想反对基督之事的人，视之为天主的敌人和灵魂的毁灭者；甚至不可\BibleQuote{向他们请安}\BibleRef{若二 1:10}，免得我们分担他们的罪恶，正如真福\Person{若望}所吩咐的。请问候与你们在一起的弟兄们。与我在一起的人也问候你们。\par

\Emph{亚历山大的司铎们}\par

7. \Signature{我，司铎\Person{柯路托}，同意以上所述，并赞同罢免\Person{亚略}及其不虔同伙的决定。}\par

\begin{itemize}
\item \Signature{司铎\Person{亚历山大}，同样}
\end{itemize}

\begin{itemize}
\item \Signature{司铎\Person{狄奥斯科若}，同样}
\end{itemize}

\begin{itemize}
\item \Signature{司铎\Person{狄奥尼修}，同样}
\end{itemize}

\begin{itemize}
\item \Signature{司铎\Person{优西比乌斯}，同样}
\end{itemize}

\begin{itemize}
\item \Signature{司铎\Person{亚历山大}，同样}
\end{itemize}

\begin{itemize}
\item \Signature{司铎\Person{尼拉辣斯}，同样}
\end{itemize}

\begin{itemize}
\item \Signature{司铎\Person{阿颇克辣提翁}，同样}
\end{itemize}

\begin{itemize}
\item \Signature{司铎\Person{阿加图斯}}
\end{itemize}

\begin{itemize}
\item \Signature{司铎\Person{内梅西乌斯}}
\end{itemize}

\begin{itemize}
\item \Signature{司铎\Person{隆古斯}}
\end{itemize}

\begin{itemize}
\item \Signature{司铎\Person{息耳瓦诺}}
\end{itemize}

\begin{itemize}
\item \Signature{司铎\Person{培罗伊斯}}
\end{itemize}

\begin{itemize}
\item \Signature{司铎\Person{阿丕斯}}
\end{itemize}

\begin{itemize}
\item \Signature{司铎\Person{普若忒里乌斯}}
\end{itemize}

\begin{itemize}
\item \Signature{司铎\Person{保禄}}
\end{itemize}

\begin{itemize}
\item \Signature{司铎\Person{居鲁士}，同样}
\end{itemize}

\begin{itemize}
\item \Emph{执事们}
\end{itemize}

\begin{itemize}
\item \Signature{执事\Person{阿摩尼乌斯}，同样}
\end{itemize}

\begin{itemize}
\item \Signature{执事\Person{玛加略}}
\end{itemize}

\begin{itemize}
\item \Signature{执事\Person{丕斯图斯}，同样}
\end{itemize}

\begin{itemize}
\item \Signature{执事\Person{亚大纳削}}
\end{itemize}

\begin{itemize}
\item \Signature{执事\Person{欧默内斯}}
\end{itemize}

\begin{itemize}
\item \Signature{执事\Person{阿颇罗尼乌斯}}
\end{itemize}

\begin{itemize}
\item \Signature{执事\Person{奥林匹乌斯}}
\end{itemize}

\begin{itemize}
\item \Person{阿弗索尼}，执事
\end{itemize}

\begin{itemize}
\item \Person{亚大纳削}，执事
\end{itemize}

\begin{itemize}
\item \Person{玛卡里}，执事，同
\end{itemize}

\begin{itemize}
\item \Person{保禄}，执事
\end{itemize}

\begin{itemize}
\item \Person{伯多禄}，执事
\end{itemize}

\begin{itemize}
\item \Person{安比提安}，执事
\end{itemize}

\begin{itemize}
\item \Person{加约}，执事，同
\end{itemize}

\begin{itemize}
\item \Person{亚历山大}，执事
\end{itemize}

\begin{itemize}
\item \Person{狄尼修}，执事
\end{itemize}

\begin{itemize}
\item \Person{阿加通}，执事
\end{itemize}

\begin{itemize}
\item \Person{波利比}，执事，同
\end{itemize}

\begin{itemize}
\item \Person{德奥纳}，执事
\end{itemize}

\begin{itemize}
\item \Person{马尔谷}，执事
\end{itemize}

\begin{itemize}
\item \Person{科莫杜斯}，执事
\end{itemize}

\begin{itemize}
\item \Person{塞拉皮翁}，执事
\end{itemize}

\begin{itemize}
\item \Person{尼隆}，执事
\end{itemize}

\begin{itemize}
\item \Person{罗曼}，执事，同
\end{itemize}

\begin{itemize}
\item \Emph{\Place{玛瑞奥提斯}的司铎。}
\end{itemize}

我，\Person{阿波罗尼}，司铎，赞同此间所写，并同意罢免\Person{阿里乌}及其不虔敬的同伙。\par

\begin{itemize}
\item \Person{英格尼}，司铎，同
\end{itemize}

\begin{itemize}
\item \Person{阿摩尼}，司铎
\end{itemize}

\begin{itemize}
\item \Person{狄奥斯库若}，司铎
\end{itemize}

\begin{itemize}
\item \Person{索斯特拉}，司铎
\end{itemize}

\begin{itemize}
\item \Person{德翁}，司铎
\end{itemize}

\begin{itemize}
\item \Person{提兰}，司铎
\end{itemize}

\begin{itemize}
\item \Person{科普雷}，司铎
\end{itemize}

\begin{itemize}
\item \Person{阿摩纳}，司铎
\end{itemize}

\begin{itemize}
\item \Person{奥里安}，司铎
\end{itemize}

\begin{itemize}
\item \Person{塞雷努斯}，司铎
\end{itemize}

\begin{itemize}
\item \Person{迪迪莫}，司铎
\end{itemize}

\begin{itemize}
\item \Person{赫拉克勒斯}，司铎
\end{itemize}

\begin{itemize}
\item \Person{博康}，司铎
\end{itemize}

\begin{itemize}
\item \Person{阿加图斯}，司铎
\end{itemize}

\begin{itemize}
\item \Person{阿基拉}，司铎
\end{itemize}

\begin{itemize}
\item \Person{保禄}，司铎
\end{itemize}

\begin{itemize}
\item \Person{塔莱莱}，司铎
\end{itemize}

\begin{itemize}
\item \Person{狄尼修}，司铎，同
\end{itemize}

\begin{itemize}
\item 执事
\end{itemize}

\begin{itemize}
\item \Person{塞拉皮翁}，执事，同
\end{itemize}

\begin{itemize}
\item \Person{尤斯图斯}，执事，同
\end{itemize}

\begin{itemize}
\item \Person{迪迪莫}，执事
\end{itemize}

\begin{itemize}
\item \Person{德默特琉}，执事
\end{itemize}

\begin{itemize}
\item \Person{毛鲁斯}，执事
\end{itemize}

\begin{itemize}
\item \Person{亚历山大}，执事
\end{itemize}

\begin{itemize}
\item \Person{马尔谷}，执事
\end{itemize}

\begin{itemize}
\item \Person{科蒙}，执事
\end{itemize}

\begin{itemize}
\item \Person{特里丰}，执事
\end{itemize}

\begin{itemize}
\item \Person{阿摩尼}，执事
\end{itemize}

\begin{itemize}
\item \Person{迪迪莫}，执事
\end{itemize}

\begin{itemize}
\item \Person{托拉里翁}，执事
\end{itemize}

\begin{itemize}
\item \Person{塞拉斯}，执事
\end{itemize}

\begin{itemize}
\item \Person{加约}，执事
\end{itemize}

\begin{itemize}
\item \Person{希埃拉克斯}，执事
\end{itemize}

\begin{itemize}
\item \Person{马尔谷}，执事
\end{itemize}

\begin{itemize}
\item \Person{德奥纳}，执事
\end{itemize}

\begin{itemize}
\item \Person{萨尔马顿}，执事
\end{itemize}

\begin{itemize}
\item \Person{卡尔蓬}，执事
\end{itemize}

\begin{itemize}
\item \Person{佐伊卢斯}，执事，同
\end{itemize}

\SourceRecord{Translated by M. Atkinson and Archibald Robertson.；Nicene and Post-Nicene Fathers, Second Series；Vol. 4.；Edited by Philip Schaff and Henry Wace.；Buffalo, NY: Christian Literature Publishing Co.,；1892.；<http://www.newadvent.org/fathers/2803.htm>.}
